\section{Co-directed aerodynamic momentum transfer}
The separately measured PZT and ZNMF fields are paired as a prospective arrangement.
The original source pairing $X_j=y_{\rm lab}-\delta$ identifies the crossing at
$X_j=11.497987$ mm and the first supported native overtaking station at $X_j=11.635200$ mm
($y_p=11.598423$ mm). One rigid placement aligns the jet centreline with the fitted plume axis
there. Both locally measured aligned droplet components initialise developed droplets;
neither the atomiser surface nor near-source velocities are relocated. Full-support
numerical fields retain their separate calibrations and masks; the report crop is unused.
Gravity is rotated with the saved frame. Interpolation requires four finite native corners
and trajectories stop at unsupported gas cells.

For a spherical, constant-diameter droplet,
$m_d=\rho_l\pi d^3/6$, $A_d=\pi d^2/4$ and
$\mathbf w=\mathbf u_g(\mathbf r_d)-\mathbf v_d$. Newton's law gives
\[
 \dot{\mathbf r}_d=\mathbf v_d,\qquad
 m_d\dot{\mathbf v}_d=\tfrac12\rho_g C_D A_d|\mathbf w|\mathbf w
 +m_d(1-\rho_g/\rho_l)\mathbf g_p.
\]
Using $Re_p=\rho_gd|\mathbf w|/\mu_g$ and $C_D=24F/Re_p$ yields the stable form
\[
 \dot{\mathbf v}_d=F(Re_p)\mathbf w/\tau_p+(1-\rho_g/\rho_l)\mathbf g_p,
 \qquad \tau_p=\rho_ld^2/(18\mu_g).
\]
The audited existing hybrid is retained: $F=1$ for $Re_p\le0.01$ and
$F=1+0.15Re_p^{0.687}$ above, with $Re_p<800$ checked at every RK stage.
The low-Re branch explicitly neglects the small finite-Re correction; zero slip has zero
drag without division by zero. No separate jet force, velocity addition or PZT-only
forcing is included. Nominal $d=10\,\mu\mathrm m$ is an aperture/Rayleigh-based estimate,
not measured sizing. Saved water/air properties at $20\,{}^\circ\mathrm C$ and assumed
101.325 kPa are retained; diameter is constant and evaporation is neglected.

Over $y_p=11.598423$--$27.335870$ mm, the reference velocity changes from 1.039291 to 2.052240 m/s
in 11.168829 ms, with transverse displacement -1.449592 mm; these values are unchanged by the audit.
The single-droplet momentum budget is
\[
 m_d[\mathbf v_d(t)-\mathbf v_d(0)]=\int\mathbf F_D\,dt+
 \int\mathbf F_{g,\rm eff}\,dt+\boldsymbol\epsilon.
\]
Axial drag impulse is $5.863643\times10^{-13}\,\mathrm{N\,s}$, effective-gravity impulse
$-5.693652\times10^{-14}\,\mathrm{N\,s}$ and momentum gain $5.294278\times10^{-13}\,\mathrm{kg\,m\,s^{-1}}$.
The RK4-stage residual is $1.243249\times10^{-27}\,\mathrm{N\,s}$; independent saved-force trapezoidal
quadrature gives $-7.544599\times10^{-18}\,\mathrm{N\,s}$. Stage-quadrature closure is expected to
roundoff and is not independent physical validation. Zero-slip, force-direction,
analytical constant-gas Stokes position/velocity and step-refinement checks pass.

Faster gas transfers momentum through drag; finite response time leaves slip.
Upward drag produces net upward acceleration only when it exceeds effective weight.
Without the imposed jet, opposing slip and gravity can decelerate droplets, but the
PZT-only experiment establishes neither stationary surrounding air nor a unique drag
history from its spatial ensemble mean. Diameter strongly affects response through
$\tau_p\propto d^2$. The seven starts cover 49.3\% of the measured entry-plane
half-maximum width and are unweighted sensitivity locations, not a whole-plume distribution.
$Stk_{D,\rm local}=\tau_p|\mathbf u_g|/D_0$ uses flow time $D_0/|\mathbf u_g|$;
$Stk_f=\tau_pf$. Local $Re_{D,c}=U_cD_0/\nu_g$ and $Sr_{D,c}=fD_0/U_c$ preserve gaps
and exclude nonpositive $U_c$, with documented $D_0=2$ mm and forcing $f=120$ Hz,
not the 10-Hz acquisition rate. These are not nozzle-exit quantities.
The prescribed gas is a steady, one-way-coupled measured mean: loading-induced gas changes
and phase-dependent motion are unresolved. Calculated paths are not measured combined-flow
results or automatically ensemble-mean droplet trajectories; no collection efficiency or
whole-plume momentum flux is established.
